\subsection{问题2：非线性物性下的耦合演化}

\subsubsection{模型求解}

在问题1中，我们将材料的各项物性参数看作常数，并据此求出了物料内部
温度和含水率随时间的变化。问题2在此基础上进一步考虑物性参数的变化，即
材料的密度、比热容、导热系数和水分扩散系数均会随着温度或含水率发生改变。
根据附录3，问题2中各物性参数的计算式为
\begin{equation}
\begin{aligned}
\rho_2(C)
&=650+128C,\\
c_{p2}(C)
&=1450+2736\frac{C}{C+1},\\
k_2(C)
&=0.21+0.38\frac{C}{C+1},\\
D_2(C,T)
&=2.4\times10^{-3}
\exp\!\left(-\frac{0.45}{C}\right)
\exp\!\left(-\frac{3850}{T}\right).
\end{aligned}
\label{eq:q2_properties}
\end{equation}

可以看出，密度、比热容和导热系数均与含水率有关，而水分扩散系数
同时受到温度和含水率的影响。因此，在物料干燥的过程中，温度和含水率发生
变化后，各项物性参数也会随之改变，进而影响下一时刻的传热和传质过程。因而，
在求解时，不能再将这些参数预先设为固定值，而需要根据每一时刻的温度和
含水率重新计算。将式\eqref{eq:q2_properties}代入式\eqref{eq:q1_system}，即可
得到问题2所需的温度场--水分场联立模型。

在前面的求解过程中，我们已经完成了控制方程和边界条件的处理。因此，对于问题2，
我们只需在每个时间步内增加物性参数的更新过程。具体求解步骤如下：

\noindent\textbf{Step1：}输入问题1中已经确定的初始条件和边界条件，并沿物料
半径方向设置201个节点。中心位置仍采用对称边界条件，物料表面采用相应的
传热和传质边界条件。

\noindent\textbf{Step2：}根据当前时刻各节点的温度$T_i$和含水率$C_i$，利用
式\eqref{eq:q2_properties}分别计算该节点处的$\rho_{2,i}$、$c_{p2,i}$、
$k_{2,i}$和$D_{2,i}$。

\noindent\textbf{Step3：}将更新后的物性参数代入传热传质方程，采用与问题1
相同的有限体积法对空间进行离散，并用BDF方法求出下一时刻各节点的温度
和含水率。若该时间步的计算结果满足收敛条件，则继续计算下一时刻；否则，
需再次更新物性参数并进行求解。

\noindent\textbf{Step4：}重复上述过程，直到计算时间达到$10800\ \mathrm{s}$。
最后输出物料在代表时刻的径向分布曲线以及前3 h内的温度场和水分场变化
情况，并对物料中心和表面的计算结果进行比较。

\subsubsection{求解结果}

使用上述方法对问题2进行计算。当$t=10800\ \mathrm{s}$时，物料中心温度为
$49.8484\ {}^\circ\mathrm{C}$，表面温度为$49.8975\ {}^\circ\mathrm{C}$，
二者仅相差$0.0491\ {}^\circ\mathrm{C}$。由此可以看出，经过3 h后，物料
中心和表面的温度已经十分接近，温度在径向上的差异较小。

在同一时刻，物料中心处的含水率为$1.7658\ \mathrm{kg/kg}$，表面处的含水率
为$1.0080\ \mathrm{kg/kg}$，二者相差$0.7578\ \mathrm{kg/kg}$。与温度的
计算结果相比，此时中心和表面的含水率仍有较大差异，说明物料内部的水分
尚未充分向表面迁移。

物料在代表时刻的温度和含水率径向分布如图\ref{fig:q2_profiles}所示。从图中
可以看出，温度曲线已经较为平缓，而含水率从中心到表面呈现出较为明显的
下降趋势，这与上述计算结果相符。

\begin{figure}[H]
\centering
\includegraphics[width=0.88\textwidth]{question2_profiles.png}
\caption{问题2在代表时刻的径向温度与含水率分布}
\label{fig:q2_profiles}
\end{figure}

前3 h内温度场和水分场的变化情况如图\ref{fig:q2_heatmaps}所示。从图可以看出，
随着加热时间的增加，物料各位置温度的差异逐渐减小；而相比之下，含水率的变化则要慢得多，直至计算结束时，中心处的含水率仍高于表面。这是由于水分扩散系数
$D_2(C,T)$会随着温度和含水率的变化而改变，使得物料内部不同位置的水分
迁移速度并不相同。因此，在前3 h内，物料的温度已基本趋于均匀，而水分仍
需较长时间才能由内部逐渐向外扩散。

\begin{figure}[H]
\centering
\includegraphics[width=0.92\textwidth]{question2_heatmaps.png}
\caption{问题2前3 h温度场与水分场的二维时空分布}
\label{fig:q2_heatmaps}
\end{figure}

综上所述，在充分考虑物料的非线性物性之后，它的温度的变化速度仍明显快于水分向外迁移的
速度。经过3 h计算，物料内部温度已经较为均匀，但中心和表面的含水率仍存在
较大差异。因此，在后续判断物料是否达到干燥要求时，不能只考虑温度变化，
还需要关注内部水分的扩散过程。
